% ===========================================================================
% SHIFT-SURV — patch for the 2026-07-28 draft
% Fixes, in order: (1) seven broken \ref/\cite, (2) figure scattering,
% (3) the figures we have but never place, (4) duplicated limitations.
% Each block says exactly what to replace.
% ===========================================================================


% ---------------------------------------------------------------------------
% [1] PREAMBLE — add these three lines after \usepackage{subcaption}
% ---------------------------------------------------------------------------
% dblfloatfix is the actual fix for "figure* lands pages later, out of order":
% stock LaTeX can only put a double-column float at the TOP of a page and never
% reorders it, which is why fig_real_res ended up inside the bibliography.
% placeins/\FloatBarrier keeps SINGLE-column floats inside their section
% (it does not constrain figure*, hence dblfloatfix as well).
\usepackage{dblfloatfix}
\usepackage{placeins}
\usepackage{flafter}    % never place a float before its first reference

% Let more floats accumulate per page instead of being deferred to the end.
\setcounter{topnumber}{3}
\setcounter{dbltopnumber}{3}
\renewcommand{\topfraction}{0.9}
\renewcommand{\dbltopfraction}{0.9}
\renewcommand{\textfraction}{0.07}
\renewcommand{\floatpagefraction}{0.75}
\renewcommand{\dblfloatpagefraction}{0.75}

% Then put \FloatBarrier immediately BEFORE each of:
%   \section{Results}, \section{Ablations and limitations},
%   \section{Discussion}, \section{Conclusion}, \section*{Acknowledgment}
% The one before \section*{Acknowledgment} is what stops a float from
% drifting into the bibliography.


% ---------------------------------------------------------------------------
% [2] BROKEN CITATION — Related Work, "Adaptation and generalization under shift"
% ---------------------------------------------------------------------------
% Currently renders as "[?], [19], [20]" on page 2: the key does not exist.
% WAS:  ...or predictions~\cite{sun2016domain,ganin2016domain,long2018conditional}.
% USE:  (the bibliography entry is sun2016deep)
%   ...or predictions~\cite{sun2016deep,ganin2016domain,long2018conditional}.


% ---------------------------------------------------------------------------
% [3] REPLACE the single fig_real_res figure* with the two split figures.
% ---------------------------------------------------------------------------
% This one edit fixes four broken refs (fig:i3lung_res x2, fig:registry_res x2)
% AND most of the scattering: two medium floats place far more easily than one
% full-width one, which is what drifted to page 9. It also finally uses
% fig_i3lung_disc/cal and fig_registry_disc/cal, which exist but are unused.
%
% DELETE the whole \begin{figure*}...\includegraphics{fig_real_res.pdf}...\end{figure*}
% block and put this one right after the I3LUNG paragraph:

\begin{figure*}[!t]
\centering
\begin{subfigure}[t]{0.49\textwidth}\centering
  \includegraphics[width=\textwidth]{fig_i3lung_disc.pdf}
  \caption{Discrimination (AUROC gap versus ERM): adaptation and domain-generalization methods.}
  \label{fig:i3lung_disc}
\end{subfigure}\hfill
\begin{subfigure}[t]{0.49\textwidth}\centering
  \includegraphics[width=\textwidth]{fig_i3lung_cal.pdf}
  \caption{Calibration (ECE gain versus ERM): adaptation, post-hoc recalibration, and generalization.}
  \label{fig:i3lung_cal}
\end{subfigure}
\caption{I3LUNG transfer results, each method versus ERM (positive is better).
No method recovers discrimination and IRM is significantly worse; recalibration
does not recover calibration. Points, random-effects pooled estimate; bars, 95\%
confidence interval.}
\label{fig:i3lung_res}
\end{figure*}

% and this one right after the registry paragraph:

\begin{figure*}[!t]
\centering
\begin{subfigure}[t]{0.49\textwidth}\centering
  \includegraphics[width=\textwidth]{fig_registry_disc.pdf}
  \caption{Discrimination (AUROC gap versus ERM); adaptation methods only.}
  \label{fig:registry_disc}
\end{subfigure}\hfill
\begin{subfigure}[t]{0.49\textwidth}\centering
  \includegraphics[width=\textwidth]{fig_registry_cal.pdf}
  \caption{Calibration (ECE gain versus ERM): adaptation and post-hoc recalibration.}
  \label{fig:registry_cal}
\end{subfigure}
\caption{Breast-cancer registry transfer results, each method versus ERM
(positive is better). Domain generalization is not applicable with two centers.
Points, pooled estimate; bars, 95\% confidence interval.}
\label{fig:registry_res}
\end{figure*}


% ---------------------------------------------------------------------------
% [4] MISSING FIGURE — fig:geometry, referenced twice in "Why alignment fails"
% ---------------------------------------------------------------------------
% Renders as "Figure ??a" / "Figure ??b" on page 7: the environment was never
% written. fig_whyfails.pdf already exists and is exactly this content
% (verified: panel a is decision angle vs AUROC gap per method; panel b is the
% covariance distance CORAL removes, annotated with the residual angle change).
% Insert immediately after the "Why alignment fails" paragraph.
%
% figure* is required, not optional: the image is 2.25:1, so at \columnwidth
% (~3.5in) the tick labels stop being legible. With dblfloatfix from [1] this
% float can also be placed at the bottom of a page, so it will not drift.

\begin{figure*}[!b]
\centering
\includegraphics[width=0.86\textwidth]{fig_whyfails.pdf}
\caption{Why alignment fails. \textbf{(a)} Angle between each method's estimated
decision direction and the target-optimal direction, against the AUROC gap, one
point per shift regime: the misalignment tracks the loss in discrimination.
\textbf{(b)} CORAL reduces the source--target covariance distance by 42 to 62
percent while leaving the decision angle essentially unchanged. The
target-optimal direction is fitted on true target labels and used only as an
analysis reference, never to fit a compared method.}
\label{fig:geometry}
\end{figure*}


% ---------------------------------------------------------------------------
% [5] MISSING FIGURES — fig:nsweep and fig:nlsweep, both in the Ablations
% ---------------------------------------------------------------------------
% Render as "Fig. ??a" and "Fig. ??" on page 7.
%
% SPACE: the draft is at 9 of 10 pages. fig:geometry + fig:nsweep + tab:meta
% already consume the remaining page, so adding fig:nlsweep as well overflows.
% Recommendation: add fig:nsweep (below), and for the nonlinear ablation DROP
% the figure and keep the sentence, whose content is a single number:
%
%   ...across the same sample-size sweep adaptation still recovers at most
%   $0.005$ AUROC, so the negative result is not an artifact of a linear
%   observation model.
%
% i.e. delete "(Fig.~\ref{fig:nlsweep})" rather than add the float. If you would
% rather keep it, it is the same shape as fig:nsweep and takes the same
% figure* treatment.
%
% CAUTION — the caption below is corrected. Panel (b) of fig5_nsweep is NOT the
% oracle headroom; it is calibration recovery (ERM minus prior-EM ECE) against
% sample size. That panel is in fact the only figure evidence for the claim,
% made in the abstract and the discussion, that label-shift recalibration
% improves with the amount of unlabeled target data, so cite it there too --
% at present the ablation text cites only panel (a) and panel (b) goes unused.

\begin{figure*}[!b]
\centering
\includegraphics[width=0.86\textwidth]{fig5_nsweep.pdf}
\caption{Neither result is a sample-size effect. Per-center sample size swept
from 250 to 2000 patients, one line per shift regime; bands, 95\% confidence
interval over ten seeds. \textbf{(a)} The best adaptation method stays within
$0.01$ AUROC of ERM at every size, and the residual gap does not close as
patients are added. \textbf{(b)} Calibration recovery by prior adjustment
(ERM minus prior-EM ECE) grows with sample size under label-prior shift, and
stays negative --- that is, harmful --- wherever the label-shift assumption is
violated.}
\label{fig:nsweep}
\end{figure*}


% ---------------------------------------------------------------------------
% [6] MISSING TABLE — tab:meta, referenced twice on page 6
% ---------------------------------------------------------------------------
% Renders as "Table ??" twice. The table existed in the July 20 draft and was
% dropped. Numbers below are taken from that draft and match the current text.

\begin{table}[!t]
\caption{Random-effects meta-analysis of the best adaptation method versus ERM,
pooled across center-to-center transfers. Positive AUROC gain favors adaptation;
both intervals exclude a clinically meaningful gain of $0.02$.}
\label{tab:meta}
\centering
\footnotesize
\begin{tabular}{@{}lccc@{}}
\toprule
\textbf{Cohort} & \textbf{Transfers} & \textbf{AUROC gain} & \textbf{95\% CI} \\
\midrule
I3LUNG (lung)            & 7 & $+0.003$ & $[-0.001, +0.007]$ \\
Public registry (breast) & 6 & $+0.005$ & $[+0.003, +0.008]$ \\
\bottomrule
\end{tabular}
\end{table}


% ---------------------------------------------------------------------------
% [7] DUPLICATED LIMITATIONS — recover roughly half a column
% ---------------------------------------------------------------------------
% The second paragraph of "Ablations and limitations" and the second paragraph
% of "Discussion" make the same four points: synthetic ground truth is a
% simplification, the real cohorts share source centers and pool correlated
% horizons, only two mechanisms are measurable on real data, single horizon.
%
% Recommendation: keep the version in "Ablations and limitations" (it sits with
% the evidence) and DELETE the limitations paragraph from the Discussion,
% replacing it with the forward-looking paragraph the supervisor asked for
% (clinical implications / multimodal), which is in snippet_discussion.tex.
% This is also the cheapest way to buy back the space the new floats need.


% ---------------------------------------------------------------------------
% [8] FIGURE STYLE — what is already consistent, and the one thing that is not
% ---------------------------------------------------------------------------
% Checked the source images directly rather than the compiled PDF.
%
% ALREADY CONSISTENT (no action):
%   - Typeface: all plots are sans-serif.
%   - Scenario colors are stable across figures: Covariate blue, Measurement
%     orange, Spurious green, Baseline hazard pink, Label-prior light blue,
%     Concept gold (fig3_dissociation, fig5_nsweep, fig_whyfails).
%   - Method-family colors are stable: Adaptation orange, Generalization teal,
%     Recalibration gold (fig_i3lung_*, fig_registry_*, fig_real_res).
%
% NOT CONSISTENT — panel lettering, three conventions in use at once:
%   (i)   fig_i3lung_disc/cal, fig_registry_disc/cal: NO letter baked in.
%         LaTeX subcaption supplies "(a)"/"(b)" below the panel. This is the
%         style the supervisor asked for, and these four are already correct.
%   (ii)  fig_whyfails: "(a)"/"(b)" baked in, centered below the panels.
%         Reads the same as (i) in print, so it is acceptable as-is, but it
%         must NOT be wrapped in subfigure or the letters will double.
%   (iii) fig5_nsweep: bold lowercase "a"/"b" baked in at the TOP-LEFT, and
%         fig_real_res additionally bakes in panel titles ("I3LUNG --
%         discrimination"). This is the convention the supervisor asked to
%         drop. fig_real_res is retired by [3]; fig5_nsweep is the only
%         remaining figure in the wrong style.
%
% CHEAPEST FIX: regenerate fig5_nsweep without the baked-in "a"/"b" and split
% it into two subfigures like [3], so every multi-panel figure in the paper
% gets its letters from LaTeX. Needs the plotting script, which is not in the
% repo. Second-best, if the script is unavailable: leave it and accept one
% figure lettered top-left, since the caption text is unambiguous either way.
